\begin{theorem}[purely analytic vanishing of the global restriction]
\label{thm:analytic-all-picard}
\Verified\ For the constructed threefold,
\[
 r=\eta^*i^*:\operatorname{Pic}(X)\longrightarrow
 \operatorname{Pic}(\mathrm{dP}_6)
\]
is the zero homomorphism.  This proof uses only holomorphic
\(\mathcal O^*\)-factors, the two finite logarithmic fillings, the cusp
argument principle, and the \(\Psi\)-equivariant toric support function.
It uses neither a global integral Leray transgression nor local integral
nearby-cycle specialization.
\end{theorem}

\begin{proof}
Let \(L\in\operatorname{Pic}(X)\).  By
\cref{thm:coinvariants}, its invariant fibre class is
\[
 E=c_1(L|_F)=kq,
 \qquad k\in\mathbb Z,
 \qquad q=uw+6\gamma\delta.
\tag{5.17-A35}
\]
The exact contractions, now used without an integral specialization
lattice, are
\[
 \iota_{v_1}q=2\eta_1,
 \qquad
 \iota_{v_2}q=-3\eta_2,
 \qquad
 [\eta_1]=3[\delta],
 \qquad
 [\eta_2]=2[\delta].
\tag{5.17-A36}
\]
Apply (5.17-A5) with \((m_1,m_2)=(3,4)\).  In the real global
coinvariants,
\[
 \bar d_1=2k[\delta],
 \qquad
 \bar d_2=-\frac32k[\delta].
\tag{5.17-A37}
\]
Equations (5.17-A4) and (5.17-A6) therefore give
\[
 0=\bar d_1+\bar d_2+\bar d_0
   =\frac{k}{2}[\delta].
\tag{5.17-A38}
\]
Since \(V_G\otimes\mathbb R=\mathbb R[\delta]\), one has \(k=0\), so
\(E=0\).  Apply \cref{lem:analytic-support-specialization} to
\(M=L|_{N_0}\).  It gives
\[
 \eta^*i^*L=\nu^*(L|_W)\simeq\mathcal O_{E_0}.
\]
This holds for every \(L\), proving \(r=0\).
\end{proof}
